\documentclass[11pt,letterpaper]{article}
\usepackage[T1]{fontenc}
\usepackage{lmodern}
\usepackage[margin=0.80in]{geometry}
\usepackage{amsmath,amssymb,booktabs,tabularx,array,longtable,microtype,xcolor}
\usepackage{fancyhdr,hyperref}
\definecolor{ink}{HTML}{17354B}
\hypersetup{colorlinks=true,linkcolor=ink,urlcolor=ink,pdftitle={CRGC W07. Nonanticipative certified opportunity assessment},pdfauthor={CRGC research program}}
\pagestyle{fancy}\fancyhf{}\setlength{\headheight}{15pt}\fancyhead[L]{\small CRGC \textbar{} Experimental research}\fancyhead[R]{\small W07 \textbar{} November 25, 2026}\fancyfoot[C]{\thepage}
\setlength{\parindent}{0pt}\setlength{\parskip}{4pt}\setlength{\emergencystretch}{2em}
\newcommand{\E}{\mathbb E}\newcommand{\R}{\mathbb R}
\begin{document}
{\small\bfseries URIEL RESEARCH \quad / \quad CRGC W07}\\[8pt]
{\huge\bfseries How much improvement survives decision-time information?}\\[7pt]
{\large Nonanticipative certified opportunity assessment and multi-call reserve viability}\\[5pt]
{\small Technical results for independent review \enspace$\cdot$\enspace November 25, 2026}

\section*{Decision summary}
On a finite storage benchmark, enforcing the information actually available to a controller ruled out a material cost reduction in six families that the analytical and perfect-foresight screens left open. The full screening hierarchy stopped in twelve of fifteen families before solving the full causal problem. Direct exact dynamic programming cost less across this small panel. The three largest gains were planted controls.

The strongest W06 three-scenario economic policy could pass every isolated reserve call and still fail the very next call. On a new 72-hour contract, the two-call standby and one-call versions failed in 20 and 13 of 48 trajectories. A budget-aware safety check passed all 48 for both price policies. The three-scenario version cost \$6.88/site-day less on average than the deterministic version. The calibrated value of an improved field controller remains unknown.

The rigorous bounds ruled out a 1\% cost improvement on one of four static AC snapshots and three of four conventional UC instances. On the one open UC case a separately registered longer run produced a feasible 1.0582\% reduction against the 30-second witness. Two stronger AC certificate jobs failed; the previously verified floors survived. Track B proves the information and viability relations conditionally, with a specific added boundedness hypothesis for a real-valued penalized infimum on an infinite domain.

\section*{Evidence contract}
The v0.6 manuscript is a separate, unchanged input. Its 26 literal source payloads produce 2,085 lines, 53 bibliography items and 91 compiled pages. All 52 W06 ID/disposition pairs match the independent ledger (42 admit, nine narrow, one hold). W06 runs here are regression tests. Each fresh track has its own freeze, sample units and objective. A and E prices are synthetic. C records benchmark cost units. AIPC remains supporting mathematical assurance. No new field observations or implementations were purchased or measured.

Every numeric finding below has a claim ID in \path{claims.csv}. Each proof, engineering premise, witness and counterexample has its own scope. The internal cross-checkers share the project specifications and compute host. The separate release receipt records the integrity and replay of the delivered files. Independent adjudication by Chat is pending.
\section{W07-B. Information, opportunity and robust viability}
\label{sec:w07formal}

\subsection{What the mathematical certificate covers}
Let $z$ be a pathwise command array, $F(z)=\sum_s w_s J_s(z_s)$ its cost, and $D$ a collection of equalities $z_{t,i}=z_{t,j}$. Each equality needs a reason. If the policy sees the same history and endogenous state on two paths at the nominal decision, its command must agree. A proposed equality between distinguishable histories can instead exclude a legitimate policy. The projection lemma explicitly requires physical feasibility of the policy's realization and equality of its observations for every pair selected by $D$. This applies to deterministic policies. A different information structure, random seed, risk constraint, or deadline requires its own coverage argument.

Write $v(S)=\inf_{z\in S} F(z)$. For a common bounded-below cost and a feasible baseline $b$ in the deadline class, two nested sets of necessary equalities give
\[
v(\mathrm{PF}) \le v(D_1)\le v(D_2)\le v(\mathrm{causal})
\le v(\mathrm{deadline})\le F(b).
\]
The maximum of validated lower floors is itself a floor. With an inclusive improvement target $\varepsilon$, a strict ceiling $F(b)-L<\varepsilon$ excludes it. At equality the target remains possible unless an independent argument excludes attainment of $L$. Original nonnegative weights preserve the bound when a missing floor is replaced by a valid fallback at its original mass. If baseline cost is missing on a positive-mass case, the weighted upper bound needs an independently justified enclosure for that cost. For example, dropping a tenth of the weight with gain 100 changes the true mean 10 to a retained-case mean of zero when every other gain is zero. These are finite-contract statements (W07-B01--B02).

\subsection{Uniform centering and a counterexample that changes the decision}
For minimization, a relaxed floor for $F+M$ is valid when $M$ has nonpositive original-weight expectation for \emph{every} covered causal policy. The Lean proof allows a different relaxed representation if it covers every admissible policy and its unpenalized cost on each projection is no larger than the original cost. For a finite tree a sufficient condition uses the history cell $h$ and an increment $m_t$:
\[
 \sum_{s:\,h_t(s)=h} w_s m_t(s,a)=0\qquad \text{for every applicable }h,a .
\]
Summing by histories and then by time yields zero for any adapted choice of actions. Zero-mass cells need no divisions. The proof also admits a checked cap $c_{t,h}$ in place of zero. If the sum of such caps over all cells and times is $\eta$, the penalized floor $L$ yields the debited floor $L-\eta$. Any deployed history-dependent state must appear in the conditioning argument. Testing only the chosen controller leaves open whether another covered controller violates the cap. This extension applies finite-sum weak duality (W07-B03--B04).

The falsification game has paths $(X,Z)=(0,0),(0,1),(1,0),(1,1)$ of respective masses $(3,1,1,3)/8$. Choose $a$ before seeing $X$, then $b$ after observing $X$ and before $Z$. Loss is $2\mathbf{1}_{a\ne X}+4\mathbf{1}_{b\ne Z}$. The constant-zero baseline costs 3. Foresight costs 0, root-only nonanticipativity costs 1 and the full adapted minimum costs 2. A feasible policy takes $a=0,b=X$. Tying $b$ across different observed values of $X$ returns the spurious floor 3 and falsely excludes the achievable gain of 1. A penalty built as conditional expected component loss less realized component loss yields the exact floor 2. Adding $4\mathbf{1}_{X=1,b=1}$ preserves zero mean on the baseline but moves the penalized minimum to 3. Its maximum causal mean defect is 2; the debited floor is 1. The nominal game enumerates 256 path arrays and eight causal policies. Across 16 additional positive-mass panels the centered floor again attains the causal value. These deliberately constructed tests demonstrate loss of validity under a false assumption. They give no generic tightness rate (W07-B05).

\paragraph{An additional real-infimum obligation.}
Pointwise integrability of every penalty does not ensure a finite floor for an infinite relaxed program. Let the original policy be $z=0$, the relaxed domain $z\le0$, the cost $F=-1$ and $M(z)=z$ in a one-scenario model. The penalty is zero on the only original policy, and every individual penalty is finite. The relaxed penalized values fill $(-\infty,-1]$. Their extended-real infimum is $-\infty$. Lean's totalized real infimum on this unbounded-below set is zero and would give a false floor against the original cost $-1$. Four separate addendum theorems establish the counterexample and prove the constructive qualification: use a verified uniform real floor, prove boundedness below for the real infimum, or use extended reals and report vacuity. W07-A's finite programs have finite cost and a feasible baseline (W07-B10).

\subsection{The order of reserve decisions}
The kernel theorem uses a finite horizon and four alternating moves: one nominal action, every admitted call, one response after that call, then every admitted disturbance. The continuation uses the public call and observed successor state. Two disturbances producing the same successor cannot inform different continuations. Induction constructs a safe contingent plan from membership in the backward kernel and proves the converse. The mathematical construction uses classical choice for witnesses and arbitrary defaults on inaccessible branches. Enlarging the call set shrinks the kernel if the older call remains admissible, the successor budget update preserves its order, the terminal sets can only shrink, and the remaining clauses of the contract are common. For a remaining count, $b\mapsto b-\mathbf{1}_{c\ne0}$ is order preserving. The sets can coincide even as the budget grows. A call-matching toy permits a safe response $u=c$ after seeing $c$ while a single response fixed before $c$ cannot match both directions (W07-B06).

\subsection{Statement correspondence, verification and prior art}
A separate script recomputed W07-A branch probabilities from the 32 released path masses in each of 15 families. It evaluated the implementation's two signed penalty increments across 465 history cells and 10,575 history--state--command combinations, including feasible actions off the baseline trajectory. Every weighted residual was exactly zero. It also checked the released ordering and penalized-floor inequality. For W07-E, pre-command available power is observable and becomes part of the current state. The automatic response fixes the new command to $q+c$. The uncertainty about later power reappears before the next command. An independent scalar evaluation of the published contract matched 230,000 admissibility bits, 18,000 pre-observation state/budget/time bits, and 250 terminal bits. Exactly 178 adjacent-budget state/time comparisons were strictly smaller. The checks assume the posted final availability window, the protected load range, the common SOC endpoint and the filed budget semantics. They address these frozen instances (W07-B07--B08).

The pinned Lean 4.30.0-rc2 build inspected 31 core and four supplementary theorem declarations, recorded 82 core and four supplementary local dependency records, and replayed four project modules using the distributed kernel checker. Nine deliberate false propositions each produced the expected rejection. Transitive axioms were confined to propositional extensionality, classical choice, and quotient soundness. The trust boundary includes the compiler and prebuilt dependencies at the nine recorded source revisions. The engineering interpreter and its correspondence checker execute outside the Lean kernel. Both local verification passes share one toolchain (W07-B09--B10).

Scenario bundling by observable information is classical in Rockafellar and Wets (1991). Weak duality for uniformly feasible penalties is a finite specialization of Brown, Smith and Sun (2010). The robust predecessor and finite target-tube interpretation are consistent with Bertsekas (1971); Horvitz (1987) addresses decisions under computational constraints. The registered contribution is the explicit contract, proof implementation, counterexamples, and finite correspondence tests. There is no established priority or deployed-field claim. The primary-source titles, links and specific comparisons are in \texttt{B/LITERATURE.md} (W07-B11).
\section{W07-A: How much improvement survives decision-time information?}

Twelve of fifteen fresh synthetic site-days admitted a certificate below the practical target before solving the full adapted problem. Six of the twelve first exclusions required tying decisions across paths that an operator could not yet distinguish. The other three site-days admitted a dispatch with a target-sized gain. Direct exact optimization was cheaper in aggregate. These conclusions concern a specified stochastic procurement contract and an ex ante feasible comparator.

\paragraph{Decision contract.}
Each site-day has six four-hour epochs. Starting with a public initial mark, five successive binary draws generate 32 positive-weight leaves and 63 decision nodes. Fifteen independent seed families were frozen across five designed regimes with weight $1/15$ each. Given current mark $z$ and epoch $t$, the published law specifies the next high-mark probability $q_{tz}/8$, protected demand, ordinary supply limit, tariff, firm backup price and wear. A forecast of next-period probability, mean price and mean ordinary capacity is issued from that law before acting. No input to the controller identifies the realized future leaf. With every observation, the controller knows the current inventory $s$, preceding signed command $r$, remaining horizon and the protected-service and endpoint contracts.

Inventory starts and ends at 2 MWh and remains in $\{0,1,2,3,4\}$. An action $a\in\{-1,0,1\}$ satisfies $|a-r|\leq1$. A charge purchases 1.25 MWh and a discharge delivers 1 MWh. All protected load is supplied. Firm backup at its higher quoted price covers ordinary supply shortages. Wear costs one dollar per MWh of internal throughput. Equal pathwise terminal stock and compulsory service remove inventory liquidation and service reductions as potential sources of gain. Unlimited contracted backup is an explicit modeling assumption.

The baseline chooses the least expensive of idle and four three-epoch (12-hour) branching rolling policies before the actual day begins. They have artificial-horizon terminal inventory slopes of 0, 11, 24 and 40 dollars/MWh, enforce reachability of the true endpoint, and optimize conditional expected cost over public branches. Every candidate is evaluated at original weights. This is a strong advantage available under an assumed known law, and its model-building and evaluation time is charged. The targeted reduction is at least $\tau=\max(15,0.01|U_0|)$ dollars/site-day. A bound excludes a family only when $U_0-L<\tau$. The initial development trial's single-slope controller cycled needlessly on its flat day. Both development panels were retained. The confirmatory freeze preceded opening the new families.

\paragraph{Certificate construction.}
The analytical tier drops intertemporal coupling. On a flat nonnegative tariff, conservation of energy with equal endpoint stocks gives the idle procurement floor because energy losses, wear and backup premia are nonnegative. Perfect foresight (PF) reveals the whole path immediately. In the next two relaxations the first one or three actions must agree across all paths sharing the corresponding observed histories. Full information about the remaining path then becomes available. Every original adapted policy projects into each relaxation at unchanged cost. The exact adapted DP recombines histories only where the tuple $(t,z,s,r)$ determines all future laws, costs, controls and endpoint obligations. An independent verifier keeps histories distinct and compares all admissible actions with both child branches at original masses.

For an alternative PF floor, define $v_{t+1}(z',s',a)$ as the minimum feasible one-epoch dispatch cost plus $11(2-s'')$ at its ending inventory $s''$, with true endpoint recoverability. After parent action $a$ reaches $s'$, penalize the realized next mark by
\[
 M_{t+1}=\sum_{j=0}^1 p(j\mid t,z)v_{t+1}(j,s',a)-v_{t+1}(z',s',a).
\]
The conditional mean is zero for every parent state and feasible action. This makes the expected penalized anticipative minimum a lower bound under a minimization objective. The separately implemented verifier checked 1,515 two-outcome centering rows and 7,845 reached Bellman rows. The zero-drift argument applies off the baseline trajectory. A valid penalty can worsen the floor. It strengthened raw PF on seven families, matched on seven and weakened it by about \$1.52 on seed 72431. It earned zero incremental first stops in the declared order. This uses established information-relaxation duality. \href{https://doi.org/10.1287/opre.1090.0796}{Brown, Smith and Sun (2010)}.

\begin{center}
\begin{tabular}{lrrl}
\hline
Frozen regime & Excluded & Demonstrated & Exact gain range (\$/site-day)\\
\hline
Flat tariff &3&0&0\\
Time of use &3&0&0 to 0.32\\
Forecast ambiguity &3&0&1.57 to 9.42\\
Rare outage &3&0&0\\
Planted delayed price spike &0&3&309 to 325\\
\hline
Total &12&3&0 to 325\\
\hline
\end{tabular}
\end{center}

The analytical test stopped three flat families, PF stopped three time-of-use families, one enforced epoch stopped one ambiguous family, and three enforced epochs stopped the remaining five. For seed 72421, the estimated maximum saving fell from \$28.12 at PF to \$22.11 after one tied epoch and \$14.52 after three. Full causal optimization achieved \$9.42, below the inclusive \$15 target. Every baseline was eligible. The original weighted mean ceiling progressed from \$227.71 analytically through \$72.35 (PF), \$70.92 (one epoch) and \$66.78 (three). The distinct penalized PF floor implied \$68.78 in isolation. The exact mean causal gain was $7988227/122880$ dollars/site-day (\$65.01), dominated by the intentionally positive regime. Eight families had zero exact gain. The predeclared feasible actions $(1,1,0,0,-1,-1)$ attained \$321, \$325 and \$309 on the three delayed-spike families. A 63-node binary tree-MILP supplied an exactly feasible causal witness at the DP objective for all fifteen families. Native dual values agreed within $10^{-6}$ dollars; their tolerance-based status is recorded separately.

\paragraph{Cost to decide.}
Each family and each arm used a fresh process. The recorded cold totals across 15 families were 1.949 s staged, 1.788 s direct DP and 13.461 s native MILP. Direct DP was quicker on nine paired launches. Baseline construction and checking took 0.133, 0.124 and 0.169 s respectively across the arms. Tier solving and checking took 0.094 and 0.249 s staged, 0.023 and 0.147 s direct, and 10.885 and 0.014 s native. Native model/solve time includes numerical-library import and model construction; its checker verifies the integer incumbent. Arm peak process RSS was approximately 21.1, 20.7 and 78.1 MiB. The complete campaign, including separate all-tier shadow runs and proof writes, took 20.13 s before the receipt. One rotated launch per arm and family is sufficient to record those costs and insufficient for a portable speed ranking. For this contract the direct exact DP is the practical choice.

\paragraph{Falsification and scope.}
A fair binary control has PF cost 5, causal cost 10 and baseline cost 25. Its improvement of exactly 15 meets the inclusive target. An artificial restriction to the baseline or a penalty centered only on its observed action falsely excludes this opportunity. On every frozen family, forecast reconstruction, repeated-prefix action checks and changed-future-suffix probes preserved the chronological information contract. A deliberately divergent same-history action failed the checker. A supplemental post-freeze reconstruction using unrecombined public histories reproduced all 75 candidate valuations and 945 selected node actions. All 480 original leaves were retained. The priced firm backup and public law remove risks that a real operator must measure. The only demonstrated material gains were designed into a delayed extreme tariff. Retire separate controller development for the twelve excluded contracts; further work on the delayed-price mechanism requires real forecast and tariff evidence and a stronger longer-horizon baseline. Relevant evidence is in \path{A/results/families.csv}, \path{A/results/records.json}, the compressed per-family proofs, \path{A/verify.py} and ledger entries W07-A01 through W07-A09.
\section{W07-E: Delivering adjacent reserve calls and recovering inventory}
\label{sec:w07e}
The W06 one-call guarantee survived reproduction. Its extension to adjacent calls failed on the archived seed and on fresh trajectories. A controller that accounts for the remaining call budget can deliver every admitted call in the finite model. On the fresh two-call panel, its three-scenario price controller outperformed the single-scenario alternative by \$6.88/site-day on average. The opportunity bounds cannot exclude a \$15/site-day improvement over the stronger controller on this panel.

\paragraph{Historical regression and counterexample (W07-E01--02).}
A fresh execution of both W06 one-call programs reproduced the stable fields of all 16 seven-/thirty-day audits, including 917,787 Bellman edge checks. Each of three policies passed 1,655 upward and 1,629 downward separate perturbations. The historical product included optional upward bids. On the seven-day trajectory of seed 69103, both continuation and three-scenario control proposed $q=1$ at hour~4 from $(s,p_{\rm prev})=(9,0)$. The first downward call applied $a=2$, reaching $(11,2)$. Independent enumeration found a single admissible nominal command at hour~5, $q=0$. The second downward call would reach SOC 12, above the operating maximum of~11.

\paragraph{Registered contract (W07-E03).}
Each of twelve fresh three-day synthetic seed families has weight $1/12$. Conditional on the budget $K$, four hidden request calendars each carry weight $1/4$. They include no requests, successive dawn down requests, successive peak up requests and mixed direction. One 72-hour window spans two midnights. State retains SOC, the last applied command, remaining budget and all announced jobs until their booked hours. Nominal and activated dispatch must serve the same facility and every booked job. Future jobs and prices enter the controller through its issued 12-hour forecast.

Power $p\in\{0,1,2\}$ is observed at the start of the hour. The final sixteen hours have public $p=2$ availability. The controller schedules $q$ before observing the operator call $c$. Its inverter applies $a=q+c$ for the entire hour. When $p>0$ and budget remains, $c$ is zero or the posted sign, $+1$ during hours 0--5 or $-1$ during hours 16--21. A delivered call consumes one budget unit. An unavailable request during an outage earns no activation payment and consumes no budget. Nominal power headroom is maintained after budget exhaustion. Each block is one stored MWh. Charging imports 1.1 grid MWh per block and discharge delivers 0.9. Inventory lies in 2--11 MWh, discharge respects the accessibility floor, feeder import lies in 0--9 MW, and the applied-command ramp limit is two blocks. Every allowed branch must permit a return to SOC~9 by hour~72. The product permits adjacent calls without an intermediate recovery deadline.

Let $V_t(b)$ be the stock and ramp states from which the remaining contract is viable. At termination it contains precisely the states at SOC~9. At earlier hours membership requires, for every allowed observed $p$, a single nominal $q$ whose physical transitions under every admitted $c$ reach $V_{t+1}(b-\mathbf{1}_{c\ne0})$. For any physically rated command, demand and grid-response extrema give $[18-18,48+22]=[0,70]$ deci-MWh inside the feeder limit $[0,90]$. Using both demand endpoints therefore preserves every eligible action even when the actual demand was observed. The scalar checker rebuilt 87,571 admitted state/action/power/budget/hour tuples. Viable state-hour counts for $b=0,1,2,3,4$ were 3,542, 3,373, 3,367, 3,365 and 3,364, excluding the common terminal layer. On six small adversarial games, direct recursive enumeration agreed at all 1,500 starting state/budget combinations. The games include uncertain outages, terminal calls, adjacent down/up, alternating signs and insufficient recovery.

\paragraph{Ceilings and realized reliability (W07-E04--06).}
W07 carries forward W06's strongest H12 price mechanism. The historical one-second fallback and optional bid market are absent from the fresh comparator. The fresh reference uses that mechanism with a common first action, three equally weighted issued-price perturbations, separate future planning and a \$65/MWh tail mark. All fresh arms use the same mandatory reserve product with no optional buyer. The standby and one-call versions assume their own limited envelope. The two robust arms apply budget-matched safety to the three-scenario and central-forecast rules. For each of 192 paths we first computed the strong reference and an exact retrospective lower cost over the robust-safety action set, revealing prices, outages and calls. Every controller that meets the full announced contract projects into this relaxation. A separate scalar checker verified all 1,847,927 integer Bellman alternatives and independently validated the applied trajectories and recorded failures. A second, issue-only sparse implementation reproduced all 49,994 completed nominations and 87 failed attempts and checked 1,633 booked jobs and 24 midnight boundaries (W07-E11). A fresh replay matched 403 deterministic files byte-for-byte. The summary matched field-for-field, with a different serialization order of policy names.

\begin{center}\small
\begin{tabular}{rrrrrrr}
\hline
$K$ & PF ceiling & Excl./48 & Standby fails & One-call fails & Robust fails$^{a}$ & Single gain$^{b}$\\
\hline
0 & 52.62 & 4 & 0 & 0 & 0 & $-4.89$\\
1 & 54.96 & 1 & 11 & 0 & 0 & $-5.56$\\
2 & 52.26 & 1 & 20 & 13 & 0 & $-6.88$\\
4 & 64.13 & 0 & 23 & 20 & 0 & $-6.78$\\
\hline
\end{tabular}
\end{center}
Ceilings and gains are in dollars per site-day. Exclusions require a ceiling strictly below the inclusive \$15 target. Each failure column has 48 paths per budget. $^a$Both robust arms have the same zero count. $^b$Three-scenario cost less single-scenario cost under the matched robust contract. At $K=0$, all four calendars have identical empty activation within each seed. These copies preserve the calendar weights and contribute one independent seed. Failed comparators have missing full-path cost and retain their weight in reliability summaries.

At $K=2$, no seed's mean gain from the single-scenario rule reached \$15/site-day. The paired seed-cluster percentile interval for its mean saving was $[-13.68,-1.33]$ using 20,000 draws. This interval describes variability among the twelve synthetic families. Complete-path ceilings spanned \$10.43--130.37/site-day. The strongest reference cost averaged \$5,396.78/site-day.

\paragraph{Activated cash, service and recovery (W07-E07).}
The primary three-scenario panel delivered 35 upward calls (31.5 grid MWh) and 36 downward calls (39.6 grid MWh). Five requests arrived at unavailable power with positive budget and 68 after budget exhaustion. The 48 complete paths served all protected load and bookings and closed at SOC~9. Across those paths, realized imports were 11,393.6 MWh costing \$796,512.30, throughput wear \$15,048.00, availability receipts \$12,576.00 and activation receipts \$248.50. Each endpoint credits \$450.00. The 156 unavailable reserve hours imply \$1,248.00 of forgone receipts under the illustrative tariff.

The recovery account pairs each activation calendar with the same seed/controller on its no-call trajectory and splits net cost differences by hour. In the $K=2$ dawn calendar, activation hours added \$30.96/site-day while subsequent noncall hours reduced cost by \$108.46/site-day. At $K=4$ on the peak calendar, the corresponding changes were $-\$104.09$ and $+\$61.62$. A call can use stock or free headroom, change subsequent feasible dispatch and affect receipts. These differences hold the synthetic exogenous path fixed.

\paragraph{Hidden-calendar diagnostic (W07-E08, exploratory).}
After primary analysis we allowed a retrospective planner to know future exogenous loads, prices and availability while revealing the request calendar only through realized calls. At each distinct call history, all indistinguishable calendars share a single action and each successor preserves its original weight. An independent checker reconstructed their observable partitions and verified 1,345,247 candidate actions and all stored attaining witnesses. The two-call mean ceiling fell from \$52.26 to \$49.50/site-day. None of the twelve seed-mean two-call bounds excludes the \$15 target. Future exogenous information remains an unresolved component of this ceiling.

\paragraph{Depth-sensitive wear and exact small-tree control (W07-E09--10).}
Each completed trajectory was revalued on its entire SOC history using the Xu et al. illustrative depth curve and two rainflow counters. Half-cycle costs were rounded to the nearest cent per depth. At $K=2$, replacing hourly wear with \$50, \$150 or \$300/kWh replacement changed the single-scenario policy's mean saving against the stronger control to $-\$6.61$, $-\$9.64$ or $-\$14.20$/site-day. These fixed-policy scores depend on an uncalibrated aging curve and the synthetic tariff. See \href{https://arxiv.org/abs/1707.04567}{Xu et al.} and \href{https://arxiv.org/abs/1703.07968}{Shi et al.}

A separate five-hour exact game uses known prices, permits two calls and assigns probability $1/2$ to an eligible call with budget remaining. Seven call paths remain. Complete action/call enumeration and a Bellman recursion retaining SOC history agree. Under linear wear the optimal expected cost was $275175/4$ cents. Its fixed policy under \$150/kWh rainflow cost $540095/8$ cents. Reoptimization lowered that to $539549/8$ cents, a difference of $273/4$ cents (\$0.6825) for this toy. Both policies chose the same initial command. Multi-day rainflow-optimal control remains open.

\paragraph{Independent contract challenge.}
A separate forward reachable-state program reconstructed all 192 retrospective minimum costs, 8,712 issued future rows from information available at issuance, and 49,994 completed command-hours. It recomputed the primary mean ceiling exactly as $8361/160$ dollars/site-day. Changing the contract from universal off-path recovery to survival on the realized path changed 135 of 192 minima. In seed 73433 with $K=2$ and no realized call, the weaker contract lowered the 72-hour minimum by \$366. Using the original kernel to price the weaker contract can omit admissible actions. No fee or dollar value was imputed to interrupted service after a failed trajectory. Fully priced service interruptions and the original W06 one-second action fallback remain outside this fresh comparison. See \path{review_checks/E_scope.md}.

\paragraph{Disposition.}
Retire extension of the one-call guarantee to consecutive activations. Retain the budget-aware safety check and stronger three-scenario economic comparator. The partial bound justifies further work on causal certification before new controller development. Availability, call probabilities, commercial qualification, service interruption and degradation remain unmeasured in the field. The primary run took 12.27 seconds from process start, including 5.08 seconds building baseline feedback tables, 0.46 applying the reference, 0.95 solving retrospective certificates, 0.76 writing them and 1.00 for independent checks. These are aggregate host timings. Reproduction, preregistration and raw cash, trajectory and failure records are under \texttt{E/}.
\section{W07-C. What remains after an unrestricted dispatch or commitment solve?}
The 1\% target refers to reducing the objective of a specific certified feasible witness on the identical input. If $u$ encloses that witness's cost from above and $l$ bounds every feasible cost from below, $l>0.99u$ rules out the inclusive target. The ratio is conservative for positive costs and positive floors. Two network sizes, paired normal and active-power-increase cases, and two UC dates with nested horizons define the entire panel. Observations are dependent within network or date. The original 48-hour UC input and initial history are preserved. Its 24-hour prefix is a separate horizon with its own endpoint.

\begin{center}\small
\begin{tabular}{lrrrr}\toprule
Network or date / hours & Witness $u$ & Cheap gap (\%) & Final gap (\%) & Cold s\\\midrule
60 & 92,695 & 2.151 & 0.229 & 1.4\\
60 stressed & 185,006 & 20.787 & 1.436 & 5.1\\
162 & 108,081 & 26.769 & 5.945 & 46.0\\
162 stressed & 120,887 & 13.801 & 4.327 & 48.5\\
04-03 / 24 & 1,203,540 & 0.495 & 0.495 & 34.6\\
04-03 / 48 & 2,065,597 & 1.614 & 1.580 & 69.3\\
10-27 / 24 & 795,360 & 1.058 & 0.958 & 35.0\\
10-27 / 48 & 1,790,662 & 0.898 & 0.898 & 41.5\\
\bottomrule\end{tabular}
\end{center}
The displayed witness cost rounds upward to whole benchmark cost units. The gap columns round for display. Decisions use exact endpoints. Cheap means a copperplate floor for AC or the unsplit continuous relaxation for UC. Time sums separate cold subprocesses, including failed higher-tier attempts, solver setup, witness recovery, certificate construction, in-run checks and serialization. Independent release audits are measured separately. Concurrent work on the same machine limits speed comparisons.

\paragraph{Static AC.} The original MATPOWER network includes reactive balance, voltage limits, capacitive charging, shunts, ideal transformer taps, oriented angle limits and apparent-power limits at both ends. A DC-initialized and a flat-start IPOPT solve each produced a center. An outward-rounded 160-bit Arb contraction encloses a zero of the original nodal balances inside a rectangular voltage/generation box. The center, free-coordinate indices, proof radius, residual and Jacobian norm are archived. Each of eight enclosures passed the remaining original constraints. A separate exact endpoint cost calculation found overlapping cost intervals between the two starts in all four snapshots.

Summing nodal balances and dropping nonnegative resistive losses and shunts gives the cheap floor when the sign assumptions hold. The stronger relaxation uses voltage products, thermal disks and convex cost epigraphs. Rational cuts and an endpoint residual correction for each numerical LP multiplier preserve weak duality over the full bounded domain. An auxiliary independent check reconstructed all 2,976 real/reactive flow rows by polarizing the two physical terminal-flow maps at 240 bits. It bounded the coefficient error over the full voltage-product boxes and verified the producer's outward row slack. The held-out tap magnitudes vary, but every phase-shift field is zero. A separate 17-degree control detected phase reversal at both terminals.

Only the normal 60-bus case excluded 1\%. Additional clique cuts reduced the stressed 60-bus ceiling from 2.063\% to 1.436\%. On both 162-bus cases the registered clique driver errored when a solver call supplied no multiplier. The primary floors remain independently checkable. A separately labeled diagnostic retained every verified prefix before the 30-second LP limit, giving 5.94476\% and 4.32681\% for the normal and stressed cases. Neither change crossed the target. The clique method samples valid positive-semidefinite directions. A full-SDP optimum remains unverified. The relaxation hierarchy follows the precedent of \href{https://arxiv.org/abs/1506.05202}{Coffrin, Hijazi and Van Hentenryck}. Inputs and licenses trace to the \href{https://github.com/power-grid-lib/pglib-opf/tree/dc6be4b2f85ca0e776952ec22cbd4c22396ea5a3}{pinned PGLib-OPF release}.

\paragraph{Conventional UC.} All four unrestricted, single-thread HiGHS solves reached their 30-second MILP limits. Their native relative gaps span 0.251\% to 1.469\%. Each incumbent commitment yielded a Z3 rational schedule. A second checker evaluated every upstream Pyomo row on the literal JSON decimals, all integrality restrictions, and the full startup and production objective. It checked 127,916 original rows. On October 27 at 24 hours, a complementary binary split was sufficient to cross 1\%. On April 3 at 48 hours, 15 nodes and eight covering leaves improved the ceiling from 1.589\% after the first split to 1.580\%. A separate checker repriced the branch multipliers on the original rows and verified that every split covers both values of a previously unfixed binary. Three of the four UC entries exclude 1\%. The registered April ceiling justified a direct search for an improving witness. Horizon, objective and global reserve constraints follow the \href{https://github.com/power-grid-lib/pglib-uc/tree/39a7f38cf4703de92f0291f0c873c2e98c789301}{pinned UC model}.

\paragraph{A feasible improvement on the open UC instance (exploratory).}
After the registered results were known, we locked a single 120-second unrestricted solve on April 3 at 48 hours. Its rational schedule reduced the objective from 2,065,596.8513 to 2,043,738.7312, an achieved reduction of 21,858.1201 benchmark cost units (1.0582\%). The independent checker evaluated the candidate against all 42,899 original constraints, integrality and original objective. This crosses the inclusive 1\% target against the fixed 30-second witness. Against the new witness, the preceding exact 15-node lower bound now limits any \emph{further} improvement to 0.5278\%. The native solve itself again reached its time limit. Building, solving and checking the new schedule took 124.10 seconds before the supplemental independent row audit. Keep this selection on the one open UC instance separate from the registered exclusion count. It demonstrates the value of longer conventional optimization under this benchmark deadline.

\paragraph{Scope and action.} The W06 regression rechecked all six earlier AC primary certificates and all four UC certificates, including 55,338 original UC rows. It carries zero weight in the fresh panel. Stop searching for a 1\% reduction on the four W07 instances with an exclusion. Use the longer conventional solve on the open UC date if its extra computation is affordable. On the three unresolved AC cases, an investment case requires a better feasible schedule or a tighter bound with defensible checking cost. These are static AC and conventional single-bus UC benchmarks. Contingencies, market settlement and field implementation costs remain open obligations.
\section{Decisions and outstanding tests}
\begin{description}
\item[Information-constrained screening.] Keep the cheap analytical tests and a direct causal DP for small exact trees. Partial nonanticipativity has value as a certificate when its cost, including baseline construction and checking, beats the available exact alternative. Re-test that trade-off at larger depth and under documented issued forecasts. A different distribution, baseline or service contract requires a new ceiling.
\item[Reserve reliability.] Use the budget-aware kernel only under its published call envelope, terminal clear window and observed availability. Before transferring it to a different product, change the amplitude contract, allow a late outage and measure service and terminal shortfalls on rejected branches. Cost comparisons of completed trajectories alone cannot establish dominance over a policy that failed the contract.
\item[Computational economics.] Advance the strong three-scenario comparator and tighter causal bounds. The finite panel supplies no observed integration-cost floor or commercial tariff. The remaining bound exceeds the target at the panel mean and supplies no promise of a realizable gain. A more expensive controller needs both a reliable feasible improvement and a costed implementation.
\item[Network and commitment.] Retire the inclusive 1\% target on the four benchmark instances already excluded. The longer unrestricted solve met the target on the open April UC case and left less than 0.528\% additional room against its improved witness. The three unresolved AC instances warrant a tighter floor or a feasible improvement before operational claims. Publish deadline failures and preserve every split leaf. A static or single-bus result cannot establish contingency performance.
\end{description}

\paragraph{Independent review.}
Use \path{W07_Independent_Review_Prompt.md} and the separate archive receipt. For every admission supply an original file, row, theorem or witness and an independently attempted falsifier. The absence of a contradiction on a finite panel leaves its external validity open. The authority for any future manuscript revision remains with the reviewer and research owner.
\end{document}
